\scanpage{4}
\mspage{2}
\unsure{Comme} \(\pi\) opère sur \deleted{\(\unsure{H^i(X_{\overline K})}\)} \(E\) par ``transport de structure'', on trouve une représentation \(\ell\)-adique (*)

\uline{Exemple} \unsure{0|2}.4. Soit \(A\) un \deleted{\unsure{schéma en|schéma}} groupe algébrique \inserted{comm.} sur \(K\). \deleted{\unsure{comm.}} alors on regarde
\[
E=T_\ell(A(\overline K))\overset{\mathrm{dfn}}{=}\varprojlim_{\nu}{}_{\ell^\nu}A(\overline K)
\]
\deleted{\unsure{lorsque|Le module}} le groupe ne change pas si on remplace \(A\) par \(\deleted{\unsure{(}}A^\circ_{\mathrm{red}}\) (NB. \(T_\ell\) est un foncteur exact à gauche nul sur les groupes finis), c'est un \deleted{groupe} \uline{module libre sur} \(\mathbb Z_\ell\) de rang égal : \(2\alpha+\deleted{\unsure{2}}\mu\), où \(\alpha\) (\(\mu\)) est la dim. de la ``partie ``abélienne'' resp. \inserted{``partie} ``torique'' de \((A_K)^\circ_{\mathrm{red}}\). Ici encore, \(\pi\) opère sur \(E\) par transport de structure.

\uline{Remarque} 0.5. On peut regarder l'exemple 0.\unsure{3|4} comme étant essentiellement \deleted{un} \inserted{le} cas particulier de l'exemple 0.4, \inserted{pour \(i=1\).} De façon précise, \inserted{on montre qu'}on a un isom canonique
\[
H^1(A_{\overline K},\mathbb Z_\ell)\simeq\deleted{\unsure{H^1(}}\operatorname{Hom}_{\mathbb Z_\ell}(T_\ell(A(\overline K)),\mathbb Z_\ell),
\]
ou ce qui revient au \(\overline{\mathrm m}\), puisque \(T_\ell(A_{\overline K})\) est libre de type fini sur \(\mathbb Z_\ell\) :
\[
T_\ell(A(\overline K))=\operatorname{Hom}_{\mathbb Z_\ell}(H^1(A_{\overline K},\mathbb Z_\ell),\mathbb Z_\ell).
\]

D'autre part, \deleted{dans 0.3, il} si \deleted{\(i=1\)} dans \unsure{0|3}.3. \(X\) est propre et \(i=1\), on trouve
\[
\begin{aligned}
H^1(X_{\overline K},\mathbb Z_\ell(1))&\simeq\varprojlim_\nu H^1(X_{\overline K},\mu_{\ell^\nu})\simeq T_\ell(\operatorname{Pic}(X_{\overline K}))\\
&\simeq T_\ell(\underline{\operatorname{Pic}}_{X/K}(\overline K))
\end{aligned}
\]
% archive sheet number: 3
